\subsection{PROPERTIES OF NORMS AND TRACES RELATIVE TO A MODULE}

If \(\mathbf{M}\) and \(\mathbf{M}'\) are two isomorphic A-modules with finite bases over \(\mathbf{K}\) , then, for all \(a \in \mathbf{A}\) ,

\[
\operatorname{Tr} _ {\mathbf {M} ^ {\prime}} (a) = \operatorname{Tr} _ {\mathbf {M}} (a), \quad \mathrm{N} _ {\mathbf {M} ^ {\prime}} (a) = \mathrm{N} _ {\mathbf {M}} (a), \quad \operatorname{Pc} _ {\mathbf {M} ^ {\prime}} (a; \mathbf {X}) = \operatorname{Pc} _ {\mathbf {M}} (a; \mathbf {X})\tag{14}
\]

for if \(f\) is an isomorphism of \(\mathbf{M}\) onto \(\mathbf{M}'\) , the matrix of \(a_{\mathbf{M}}\) with respect to a basis \(\mathbf{B}\) of \(\mathbf{M}\) over \(\mathbf{K}\) is the same as the matrix of \(a_{\mathbf{M}'}\) with respect to the basis \(f(\mathbf{B})\) of \(\mathbf{M}'\) .

PROPOSITION 1. Let \(\mathbf{M} = \mathbf{M}_0 \supset \mathbf{M}_1 \supset \cdots \supset \mathbf{M}_r = \{0\}\) be a decreasing sequence of submodules of an A-module \(\mathbf{M}\) such that each of the K-modules \(\mathbf{P}_i = \mathbf{M}_{i-1} / \mathbf{M}_i\) ( \(1 \leqslant i \leqslant r\) ) admits a finite basis. Then the K-module \(\mathbf{M}\) admits a finite basis and

\[
\operatorname{Tr} _ {\mathbf {M}} (a) = \sum_ {i = 1} ^ {r} \operatorname{Tr} _ {\mathrm{P} _ {i}} (a), \quad \mathrm{N} _ {\mathbf {M}} (a) = \prod_ {i = 1} ^ {r} \mathrm{N} _ {\mathrm{P} _ {i}} (a)\tag{15}
\]

\[
\mathrm{Pc} _ {\mathbf {M}} (a; \mathbf {X}) = \prod_ {i = 1} ^ {r} \mathrm{Pc} _ {\mathbf {P} _ {i}} (a; \mathbf {X}).
\]

Let \(B_{i}^{\prime}\) be a basis of \(P_{i}\) over K; then a system of representatives \(B_{i}\) of \(B_{i}^{\prime}\) (mod. \(M_{i}\) ) is a basis of a supplementary submodule of the K-module \(M_{i}\) in the K-module \(M_{i-1}\) (II, § 1, no. 11, Proposition 21). The union B of the \(B_{i}\) ( \(1 \leqslant i \leqslant r\) ) is a basis of M over K. Let \(X_{ii}\) be the matrix of the endomorphism \(a_{P_{i}}\) with respect to the basis \(B_{i}^{\prime}\) . It is immediate that the matrix of \(a_{M}\) with respect to B is of the form

\[
\left( \begin{array}{c c c c c} X _ {r r} & X _ {r, r - 1} & \ldots & X _ {r, 1} \\ 0 & X _ {r - 1, r - 1} & \ldots & X _ {r - 1, 1} \\ \cdot & \cdot & \cdot & \cdot & \cdot \\ 0 & 0 & \ldots & X _ {1 1} \end{array} \right)
\]

and the proposition follows from formulae (4), (5) and (6) of no. 1 and formula (31) of § 8, no. 6.

PROPOSITION 2. Let A, A' be two K-algebras, M an A-module and M' an A'-module. Suppose that M and M' are free K-modules of respective dimensions n and n' and consider M ⊗K M' as an (A ⊗K A')-module (§ 4, no. 3, Example 2). Then, for a \ensuremath{\in} A and a' \ensuremath{\in} A',

\begin{enumerate}
\def\labelenumi{(\arabic{enumi})}
\setcounter{enumi}{15}
\tightlist
\item
\end{enumerate}

\[
\operatorname{Tr} _ {\mathbf {M} \otimes \mathbf {M}} (a \otimes a ^ {\prime}) = \operatorname{Tr} _ {\mathbf {M}} (a) \operatorname{Tr} _ {\mathbf {M} ^ {\prime}} (a ^ {\prime})\tag{17}
\]

\[
\mathrm{N} _ {\mathbf {M} \otimes \mathbf {M} ^ {\prime}} (a \otimes a ^ {\prime}) = (\mathrm{N} _ {\mathbf {M}} (a)) ^ {n ^ {\prime}} (\mathrm{N} _ {\mathbf {M} ^ {\prime}} (a ^ {\prime})) ^ {n}.
\]

Formula (16) follows from II, § 4, no. 4, formula (26) and formula (17) from § 8, no. 6, formula (33).

